%! TEX root = esai-en.tex
% the main TeX file which is intended to compile, :VimtexReload after adjustment
% :h vimtex-tex-root
\documentclass[../main.tex]{subfiles}
%ensure the class of docs

\begin{document}

\subsection{Developed Translation}%
\label{sub:Developed Translation}

% I was born and raised in the city of Parepare. When I was a kid, I would often go for recreation around the city by walking. I used to walk with my friends from our home to the recreation place. While we were walking, we saw the city scape. The city was good but not as good as it now.

I was born and raised in Parepare city. During my childhood, I would often go for recreation around the city by walking.
I would walk with my friends from our house to the recreational areas. As we strolled, we admired the cityscape. The city was pleasant, although not as developed as it is today.

The city's recent development can be attributed to its strategic location and the support it received from the third president of Indonesia during his lifetime, Eyang Habibie.
This development has led to the initiation of major infrastructure projects such as the construction of a railway and a football stadium. The railway project is currently in its initial phase but is intended to eventually connect the capital of Makassar to Parepare. Meanwhile, the stadium has undergone revitalization in recent years.

The urban development that has occurred turns out to be growing in line with the population of Parepare city. Based on data from the Central Statistics Agency (BPS) in 2021, the population of Parepare increased by 5\% from 2019. The increasing population and development may have negative impacts on the city, as stated by most researchers, for example, urban sprawl.

As an architecture lecture, I propose a number of solutions to anticipate those impact.
First, enhancing the quality of architecture or urban planning education. Education becomes the main key for architecture students as they enter the field.

As an architecture lecture, I propose a number of solutions to anticipate those impact.
Firstly, enhancing the quality of architecture or urban planning education to create more realiable architect.
Secondly, embracing the concept of sustainable development to promote social, economic, and technological progress with efficient use of natural resources and energy.





\section{Open Ai Translation}%
\label{sec:Open Ai Translation}

I was born and raised in a city called Parepare.
This city is one of the cities in Indonesia that never escapes the attention of many people.
People know it as a religious, tourism, trade, and education city. Not a few people also know that this city is the birthplace of Eyang Habibie, the third president of the Republic of Indonesia. This popularity has driven the city to continue to advance and develop.

The current development of the city of Parepare shows a significant increase. At the end of 2022, Parepare inaugurated the construction of the first phase of a railway connecting Parepare to Makassar (the capital of South Sulawesi). In addition, the city has successfully made the BJ. Habibie stadium the home base of the Persatuan Sepakbola Makassar (PSM) football club. This development, without a doubt, is the foundation of Parepare city as a magnet for thousands of people.

Based on data from the Central Statistics Agency (BPS) in 2021, the population of Parepare increased by 5\% from 2019.
This data indicates the occurrence of urbanization that has the potential to have a negative impact on the city's future. One of the negative impacts is the disorderly arrangement of the city.


Therefore, it is important to formulate efforts to address these impacts.
In the field of architecture, these efforts are, first, to improve the quality of urban-related education. This starts at the university level. If the quality of education at universities is good, the ability of architecture graduates to overcome the impacts of urbanization in their areas can increase. This can be realized through improving the quality of architecture lecturers at universities, especially in the regions.

Second, designing by considering the concept of sustainability that focuses on the development of cities and settlements that are inclusive, safe, resilient, and sustainable. This concept will provide the city with a strong foundation for urban problems, especially urbanization issues.
There are several aspects to this concept, namely economic, energy, ecological, participation, and equality aspects. Several cities have implemented these aspects, one of which is the city of Jepara. This city has implemented at least three aspects, namely economic, energy, and participation aspects.

Designing with a focus on sustainability involves creating inclusive, safe, resilient, and sustainable cities and settlements. This approach addresses urban challenges related to urbanization and encompasses economic, energy, ecological, participation, and equality considerations. Cities like Jepara have successfully implemented these aspects, including economic, energy, and participation factors.

Secondly, designing with a focus on sustainability involves creating cities and settlements that are inclusive, safe, resilient, and sustainable.
This approach establishes a solid foundation for addressing urban challenges, particularly those related to urbanization.
The concept encompasses various aspects, including economic, energy, ecological, participation, and equality considerations.
Some cities, such as Jepara, have successfully implemented these aspects. Jepara, for instance, has incorporated at least three key aspects: economic, energy, and participation.

The implementation of these three aspects is as follows: in terms of the economic aspect, the city of Jepara builds a gathering place for business activities. This place will allow people to compete healthily, share knowledge, and carry out various activities in urban areas on a daily basis. In terms of energy, the city uses solar panels as an energy source in many places. In terms of participation, the city empowers the community to mitigate the impacts of the TJB Steam Power Plant in the northern part of Jepara. This implementation will certainly support the sustainability not only of the city of Jepara but can also be realized in other cities, thus creating the sustainability of a city as a whole.

I plan to continue my education to the doctoral level because this level provides an opportunity to sharpen sensitivity to human and environmental issues. This sensitivity will allow me to contribute to addressing city problems such as urbanization in the future. For example, continuing to teach at my alma mater with sensitivity to deeper surrounding issues.

During my studies, I will take the course-based study track. I will focus on several courses related to public space design, urban development management, and sustainable architecture in my studies. Furthermore, if my dissertation proposal is accepted, I will develop a dissertation title on the identity of coastal public spaces. The reason is that local identities are starting to erode and change. Also, the standard and uniform characters of a place are increasingly developing, largely influenced by urbanization.

I plan to continue my studies at a domestic university. The reason is that domestic universities now have very good quality. For example, there are three architecture study programs from domestic universities that are included in the ranks of the top 1000 best architecture study programs in the world. This achievement is certainly the result of a process that consumes a lot of time and energy.

When I studied for my Master's degree at Diponegoro University, I felt a spirited learning atmosphere. On one occasion, a lecturer asked me about the Bugis people's food. He asked, "People there (Bugis tribe) don't eat fish, right?" My response was quick, "They eat fish, Prof." At that moment, he confused the class, especially me, then countered, "They don't eat fish there, they only eat \textit{ikang}." I immediately nodded and chuckled, followed by other classmates. In addition to the learning atmosphere, I assessed the very good quality of knowledge. Most of the lecturers who taught me were doctors, and not a few were professors. Sometimes there were also foreign lecturers who took the time to teach as part of the Visiting Professor program. My experience is an illustration that domestic universities continue to develop for the better.

Now, domestic universities have produced prominent doctoral alumni who have made significant contributions to the country. Some of them are Moh. Mahfud MD, Agung Gumiwang Kartasasmita, and Ahmad Basarah. In addition to doctoral degree holders, many national figures have received education in Indonesia, both bachelor's and master's degrees. One of them, Mr. Jusuf Kalla, is the 10th and 12th Vice President. He is well known internationally as a peace figure. Jusuf Kalla, who is now the chairman of the Indonesian Red Cross (PMI), has provided a lot of meaning for Indonesia. These figures indicate that domestic graduates have their own value and preparation to serve the country.

The country is like a blank canvas that requires strokes of various colors. Those brushstrokes are the contributions of the nation's children that have been present since the Dutch East Indies era until now. In the past, these contributions were generally bloodshed. Now, those contributions can be very diverse. For example, writings, designs, applications, movements, and so on. The strokes on that canvas will form a beautiful painting that is enjoyed by everyone.

When I was in high school, I was active in school organizations, especially the Red Cross. At that time, we held organizational activities. The most memorable activities were, first, when we held a blood donation drive at school attended by several teachers and students. Second, when we organized training sessions for cadres such as first aid, the use of common medical equipment, and search and rescue.

During my undergraduate studies, I joined the architecture student organization outside of campus, namely PAMIY (Panguyuban Mahasiswa Arsitektur Yogyakarta). At that time, we organized several activities. First, we designed and built a small library for the mosque in the Parangritis village, Yogyakarta. Second, we held sports and arts competitions among architecture students in Yogyakarta. In addition to being active in organizations, I once served as a teaching assistant for the Detail Engineering Design course. I helped students apply the theories just explained by the lecturer. Sometimes there were students who needed further guidance, so I usually made time to teach them more intensively to catch up with other classmates.

During my Master's studies, I tended to be more active in writing than participating in extracurricular activities. The reason is the spread of a virus that hindered activities as usual. Nevertheless, it allowed me to achieve a GPA of 4.00 and the title of \textit{Cum laude}. In addition, I was able to publish an article in a reputable national journal. My writings, including my thesis, covered public spaces that I care about deeply because they are related to human life, including social, health, and even welfare aspects.

In the third semester of my Master's studies, I briefly worked on a project to renovate a madrasah school while completing my thesis proposal. This renovation required the addition of inclusive design and a garden. So, as part of the drafting team, I designed a ramp as requested during the site visit at the initial project meeting. Additionally, we also added a garden in front of the classrooms. This will certainly result in a universal design, meaning its use encompasses everyone in general, without physical, age, and gender limitations.

A few months after completing my Master's degree, I worked as a lecturer at Muhammadiyah University Parepare until now. This university has just opened the Regional and Urban Planning study program. In my work as a lecturer, I have the opportunity to teach two courses. In addition, I am also actively involved in institution activities such as scientific paper writing workshops for young lecturers.

Once I have completed my doctoral studies, I will continue my work with better skills and insights. One of the main tasks of a lecturer is to carry out the Tri Dharma of higher education. Some of them are, first, teaching urban studies in the undergraduate program. Second, conducting research that can be published in national or international journals. Lastly, conducting community service in the Ajatappareng area (several areas around Parepare) to apply the results of teaching and research.

All my contributions are inseparable from the improvement of personality aspects, namely strengths, weaknesses, and experiences that I have. First, I feel I have strengths such as a strong work ethic, a high spirit to learn new things, and strong perseverance. Second, I have weaknesses such as being impatient, lacking detail in completing tasks, and perfectionism. Third, I have proud experiences such as graduating \textit{Cum laude} during my Master's, working on the construction of a madrasah school, and being appointed as a civil servant lecturer at the Bacharuddin Jusuf Habibie Institute of Technology (ITH). Lastly, my less proud experience is failing to obtain LPDP and BPI scholarships before.

In addition, the contributions I will make to the university after completing my doctoral studies are tied to the independent learning program. Independent learning is a program that emphasizes freedom for students. Students can take courses outside their study programs and engage in activities beyond the campus scope. This allows lecturers to easily develop learning focused on academic outcomes rather than binding rules. Improving personality aspects and implementing adaptive programs are likely to maximize my contributions.

These contributions aim to enhance learning so that the Golden Indonesia mission can be achieved by 2045. Therefore, not only one individual, but cooperation from various parties will greatly help in achieving this goal. This cooperation includes, first, the role of parents in shaping children's characters from an early age. Second, the role of the community in shaping constructive values. Lastly, the role of the government in managing and facilitating learning.

Excellent education will create a generation that will bring a brilliant future for the nation. For example, good architectural or urban education will encourage generations to develop cities that are friendly to the impacts of urbanization. Education can be realized not only by changing policies but by starting with a change in paradigm and requiring a long process.

\subsection{proofread AI}%
\label{sub:proofread AI}

Act as a proofreader and review the given text. Feel free to rephrase sentences or make changes to enhance clarity but maintain the overall tone and style of the original. Show all changes in bold so I can see what has been updated.

I was born and raised in a city called Parepare. This city is **one of the notable cities in Indonesia that attracts** the attention of many people. People know it as a religious, tourism, trade, and education city. Not a few people also know that this city is the birthplace of Eyang Habibie, the third president of the Republic of Indonesia. This popularity has driven the city to continue to advance and develop.

The current development of Parepare is experiencing a significant increase. At the end of 2022, Parepare inaugurated the construction of the first phase of a railway connecting Parepare to Makassar, the capital of South Sulawesi. Additionally, the city has successfully designated the BJ. Habibie Stadium as the home base of the Persatuan Sepakbola Makassar (PSM) football club. Undoubtedly, this development lays the foundation for Parepare city to become a magnet for thousands of people.

Based on data from the Central Statistics Agency (BPS) in 2021, the population of Parepare increased by 5 compared to 2019. This data suggests urbanization is happening, which could lead to negative consequences for the city. One such consequence is the disorderly development of the city.

Therefore, it is crucial to formulate strategies to mitigate these impacts. In the field of architecture, these strategies primarily involve enhancing the quality of urban-focused education, starting at the university level. **When** the quality of education at universities is high, architecture graduates' capacity to address the challenges of urbanization in their respective areas can be significantly improved. This can be achieved by enhancing the quality of architecture professors at universities, particularly in regional settings. Another key strategy is to incorporate sustainability principles into design, emphasizing the development of cities and settlements that are inclusive, safe, resilient, and sustainable. This approach will establish a solid groundwork for addressing urban issues, particularly those related to urbanization. This sustainability concept encompasses various aspects, including economic, energy, ecological, participatory, and equality considerations. Numerous cities have already integrated these aspects into their planning, with Jepara being a notable example. Jepara has successfully implemented at least three key aspects: economic, energy, and participatory elements.

---
Act as a proofreader and review the following text. Feel free to rephrase sentences or make changes to enhance clarity but maintain the overall tone and style of the original. Show all changes in bold so I can see what has been updated.

The economic aspect is implemented by the city of Jepara through the construction of a business hub. This hub will promote healthy competition, knowledge sharing, and various urban activities on a daily basis. In terms of energy, the city utilizes solar panels as a primary energy source across multiple locations. Regarding participation, the city encourages community involvement to mitigate the impacts of the TJB Steam Power Plant in the northern region of Jepara. This initiative will significantly bolster the sustainability of not only Jepara but also other cities, thereby fostering overall urban sustainability.

I aspire to pursue doctoral studies to further refine my awareness of human and environmental concerns. This heightened awareness will enable me to contribute towards addressing urban challenges, like urbanization, in the future. For instance, I aim to continue teaching at my alma mater while remaining attuned to broader contextual issues.

When I pursued my Master's degree at Diponegoro University, I experienced a vibrant learning environment. During one class, a lecturer inquired about the traditional food of the Bugis people. He asked, "The Bugis tribe doesn't consume fish, correct?" I promptly replied, "They do eat fish, Prof." This response puzzled the class, particularly me. The lecturer then clarified, "In that region, they do not eat fish, they solely consume **\textit{ikang}**." I nodded in agreement and chuckled, as did my classmates. Along with the enriching learning atmosphere, I appreciated the exceptional quality of education. The majority of my instructors held doctoral degrees, and several were professors. Occasionally, foreign professors participated in the Visiting Professor program, dedicating their time to teaching. My personal encounter exemplifies the continuous advancement of local universities for the better.

Now, domestic universities have produced prominent doctoral alumni who have made significant contributions to the country. Some of them are Moh. Mahfud MD, Agung Gumiwang Kartasasmita, and Ahmad Basarah. In addition to doctoral degree holders, many national figures have received education in Indonesia, earning both bachelor's and master's degrees. One of them, Mr. Jusuf Kalla, is the 10th and 12th Vice President. He is well known internationally as a peace figure. Jusuf Kalla, who is now the chairman of the Indonesian Red Cross (PMI), has provided a lot of meaning for Indonesia. These figures indicate that domestic graduates have their own value and preparation to serve the country.

The country is like a blank canvas that requires strokes of various colors. Those brushstrokes are the contributions of the nation's children that have been present since the Dutch East Indies era until now. In the past, these contributions were generally **marked by** bloodshed. Now, those contributions can be very diverse. For example, writings, designs, applications, movements, and so on. The strokes on that canvas will form a beautiful painting that is enjoyed by everyone.

---
When I was in high school, I was actively engaged in school clubs, particularly the Red Cross. During that period, we conducted diverse organizational events. The most remarkable activities included, firstly, hosting a blood donation campaign at school that was joined by several teachers and students. Secondly, we arranged training workshops for members on topics like first aid, operating basic medical tools, and search and rescue techniques.

During my undergraduate studies, I was a member of the architecture student organization PAMIY (Panguyuban Mahasiswa Arsitektur Yogyakarta) outside of campus. We organized various activities. Firstly, we designed and constructed a small library for the mosque in the Parangritis village, Yogyakarta. Secondly, we hosted sports and arts competitions for architecture students in Yogyakarta. Apart from my involvement in organizations, I served as a teaching assistant for the Detail Engineering Design course. I assisted students in applying the theories discussed by the lecturer. Occasionally, students required additional support, so I dedicated time to provide them with more intensive guidance to help them catch up with their peers.

During my Master's studies, I was more focused on writing than engaging in extracurricular activities due to the disruption caused by a virus. This situation enabled me to achieve a GPA of 4.00 and earn the distinction of **\textit{Cum laude}**. Moreover, I successfully published an article in a prestigious national journal. My written works, which encompassed my thesis, delved into public spaces that hold significant importance to me as they are intertwined with various aspects of human life, such as social, health, and welfare.

---

In the third semester of my Master's studies, I briefly participated in a project to enhance a madrasah school while concurrently developing my thesis proposal. This enhancement involved incorporating inclusive design elements and establishing a garden. As a member of the design team, I was tasked with creating a ramp following a request made during the initial project meeting on-site. Furthermore, we introduced a garden in front of the classrooms. This initiative will undoubtedly lead to a universal design, ensuring that it caters to a wide range of users without limitations based on physical abilities, age, or gender.

A few months after obtaining my Master's degree, I began working as a lecturer at Muhammadiyah University Parepare, where I am currently employed. The university recently introduced the Regional and Urban Planning program. As a lecturer, I teach two courses and actively participate in institutional activities, including organizing workshops on scientific paper writing for junior faculty members.

Once I complete my doctoral studies, I will enhance my skills and insights for continued work. A key responsibility of a lecturer is to fulfill the Tri Dharma of higher education. This includes teaching urban studies in the undergraduate program, **engaging in research suitable for publication in national or international journals**, and providing community service in the Ajatappareng area (various locations around Parepare) to implement the outcomes of teaching and research.

All my contributions are closely tied to enhancing aspects of my personality, including strengths, weaknesses, and experiences. Firstly, I recognize strengths such as a robust work ethic, a keen eagerness to acquire new knowledge, and strong perseverance. Secondly, I acknowledge weaknesses like impatience, a tendency to overlook details in task completion, and a struggle with perfectionism. Thirdly, I take pride in experiences such as graduating Cum Laude in my Master's program, contributing to the construction of a madrasah school, and serving as a civil servant lecturer at the Bacharuddin Jusuf Habibie Institute of Technology (ITH). Lastly, I reflect on a less proud experience of not securing the LPDP and BPI scholarships in the past.

In addition, the contributions I will make to the university after completing my doctoral studies are connected to the independent learning program. Independent learning is a program that highlights student freedom. Students have the opportunity to enroll in courses beyond their academic programs and participate in activities that extend beyond the campus boundaries. This flexibility enables instructors to create learning experiences that prioritize academic achievements over rigid regulations. Enhancing personal development and implementing adaptable initiatives are expected to amplify my impact.

These contributions are intended to improve learning to support the Golden Indonesia mission by 2045. Therefore, collaboration from multiple parties will be crucial in reaching this objective. This collaboration involves, firstly, the responsibility of parents in molding children's characters from a young age. Secondly, the involvement of the community in instilling positive values. Lastly, the government's role in overseeing and enabling the learning process.

Excellent education will foster a generation that can pave the way for a prosperous future for the nation. For instance, a strong foundation in architecture or urban planning can inspire future generations to build cities that can effectively address the challenges of urbanization. Education should not only involve policy adjustments but also necessitate a shift in mindset, which demands a sustained effort over time.















Nowadays, it shows an increasing
Nowadayas, it is a popular destination for sightseeing, study and business.

I was born and raise in Parepare city. This city is popular in South Sulawesi Island because





\bibliographystyle{apalike}
% \bibliography{../filename.bib} % required for citecompletion, not work in mainfile
\end{document}
